\clearpage
\setcounter{page}{1}
\maketitlesupplementary

\section{Implementation Details}
\label{sec:suppl_impl}

\paragraph{Training augmentation.}
Every image passes through random horizontal flipping and then one of two
paths: a resize to a randomly sampled scale, or a stronger sequence of
resize, random-size crop, and multi-scale resize, matching the
Grounding DINO recipe~\cite{groundingdino}. At inference all augmentations are
disabled and evaluation follows COCO metrics.

\paragraph{Support-side configuration.}
FFCP uses a fixed context rank across domains, and the Hard update radius, the
routing temperature, and the Soft authority budget are shared across the six
benchmarks. The class-wise estimate is shrunk toward a cross-category drift
subspace in the 1-shot regime and receives progressively more weight as the
number of support instances grows.

\section{Additional Quantitative Analysis}
\label{sec:suppl_extra}

\paragraph{Hyperparameter sensitivity.}
Fig.~\ref{fig:hyperparam} in the main paper summarizes the sensitivity of the
five shared hyper-parameters on 5-shot Clipart1k; the complete per-value sweep
tables are provided in the released code. Each sweep varies one quantity while
keeping the support split, checkpoint, and optimization budget fixed, and the
shared cross-domain configuration is selected from the central stable region
rather than tuned per domain.

\section{Additional Qualitative Analysis}
\label{sec:suppl_qual}

\paragraph{Cross-domain qualitative comparison.}
Fig.~\ref{fig:qualitative_six_domains} shows one representative query per
target domain, with the ground truth, the text-only baseline, and the full
model as columns. Both methods share the support split, checkpoint,
preprocessing, score threshold, and post-processing.

\begin{figure*}[t]
    \centering
    \includegraphics[width=0.72\textwidth]{fig/qualitative_six_domains.png}
    \caption{Qualitative comparison across six CD-FSOD target domains in the
    5-shot setting. Each row is one domain; columns are the ground truth, the
    text-only Grounding DINO baseline, and the full S2H-CD-FSOD model.}
    \label{fig:qualitative_six_domains}
\end{figure*}

\paragraph{FFCP interventions.}
Fig.~\ref{fig:intervention_examples} visualizes the foreground-frozen
interventions that produce the context-sensitive directions used by CHSD. Each
ground-truth box is accompanied by jittered regions at increasing overlap with
the object, which retain realistic partial-object and boundary cues while
changing the surrounding context. These views expose the neighboring objects
and domain-specific background patterns that drive false positives, and explain
why the context-sensitive subspace is removed from the identity anchor.

\begin{figure*}[!t]
    \centering
    \includegraphics[width=0.6\textwidth]{fig/intervention_examples.png}
    \caption{Foreground-frozen interventions across six target domains. Each row
    shows the protected ground-truth box and several jittered regions that change
    only the surrounding context. These probes identify where external context
    perturbs the otherwise unchanged target representation.}
    \label{fig:intervention_examples}
\end{figure*}

\paragraph{Failure cases.}
Fig.~\ref{fig:failure_cases} retains representative failures so that the
qualitative analysis is not reduced to successful examples: severe occlusion,
very small objects, atypical support prototypes, and residual false positives
in the full model. Authority routing mitigates support-induced errors but does
not remove the intrinsic ambiguity of an underspecified support set.

\begin{figure*}[t]
    \centering
    \includegraphics[width=\textwidth]{fig/failure_cases.png}
    \caption{Representative failure cases. Red boxes mark false positives,
    yellow dashed boxes mark missed ground truths, and purple labels mark
    incorrect category assignments.}
    \label{fig:failure_cases}
\end{figure*}


